% GENERATED FILE. Edit canonical lessons, then run npm run generate:books.
% canonical-source-hash: fnv1a64:f1157d1c36956e52
% canonical-lessons: FR-C355-licencier, FR-C355-licenciement, FR-C355-demission, FR-C355-demissionner, FR-C355-preavis

\chapter{Quitter son travail}
\label{ch:fr-quitter-son-travail}
\hlchaptermodality{355}

\begin{chapteropening}
\textbf{By the end of this chapter:} \emph{I can talk about leaving a job: being laid off, a dismissal, resigning, a resignation and the notice period.}

Complete the last lesson of chapter 355: I can talk about leaving a job: being laid off, a dismissal, resigning, a resignation and the notice period.
\end{chapteropening}

\section[licencier]{licencier — to dismiss, to lay off}

\label{lesson:FR-C355-licencier}



\begin{quote}
Before the new one: say the French for to apply, then the French for the hiring.
\end{quote}

\subsection*{You'll want to know: licencier}
\textbf{licencier} — \textquotedblleft{}to dismiss, to lay off\textquotedblright{}.

\begin{quote}
L'usine a licencié cinquante ouvriers.
\end{quote}

\begin{quote}
\emph{The factory laid off fifty workers.}
\end{quote}

Formal; in speech people also say \emph{renvoyer}.

\begin{cousinweb}
\emph{licencier} is built on \emph{licence} \textquotedblleft{}leave, permission\textquotedblright{}, from Latin \emph{licentia} \textquotedblleft{}freedom, licence\textquotedblright{} (Lewis \& Short), from \emph{licēre} \textquotedblleft{}to be allowed\textquotedblright{}: first, to give someone leave to go.
\end{cousinweb}

\begin{grammarlens}[title={What you've built}]
One more word for telling, explaining and describing at B1.
\end{grammarlens}

\subsection*{Your turn}
\begin{itemize}
\raggedright
  \item \emph{Say it:} \emph{licencier}
  \item \emph{Say it:} \emph{licencier}, once more
  \item \emph{Say it:} a sentence of your own with \emph{licencier}
  \item \emph{Recall:} say the French for to apply, then the French for the hiring, then say \emph{licencier} again
\end{itemize}

\begin{tcolorbox}[breakable,colback=teal!4,colframe=teal!35!black,title={Before you move on}]
What does licencier mean? (\textquotedblleft{}To dismiss, to lay off\textquotedblright{}.) Say it once more.
\end{tcolorbox}
\section[le licenciement]{le licenciement — the dismissal, the layoff}

\label{lesson:FR-C355-licenciement}



\begin{quote}
Before the new one: say the French for the hiring, then the French for to dismiss.
\end{quote}

\subsection*{You'll want to know: le licenciement}
\textbf{le licenciement} — \textquotedblleft{}the dismissal, the layoff\textquotedblright{}.

\begin{quote}
Il a reçu sa lettre de licenciement.
\end{quote}

\begin{quote}
\emph{He received his letter of dismissal.}
\end{quote}

\emph{Un licenciement économique} — a layoff for economic reasons.

\begin{cousinweb}
\emph{licenciement} is built from \emph{licencier} and the suffix \emph{-ment}. The verb rests on \emph{licence}, from Latin \emph{licentia} \textquotedblleft{}freedom, licence\textquotedblright{}, from \emph{licēre} \textquotedblleft{}to be allowed\textquotedblright{} (Lewis \& Short).
\end{cousinweb}

\begin{grammarlens}[title={What you've built}]
One more word for telling, explaining and describing at B1.
\end{grammarlens}

\subsection*{Your turn}
\begin{itemize}
\raggedright
  \item \emph{Say it:} \emph{le licenciement}
  \item \emph{Say it:} \emph{le licenciement}, once more
  \item \emph{Say it:} a sentence of your own with \emph{le licenciement}
  \item \emph{Recall:} say the French for the hiring, then the French for to dismiss, then say \emph{le licenciement} again
\end{itemize}

\begin{tcolorbox}[breakable,colback=teal!4,colframe=teal!35!black,title={Before you move on}]
What does le licenciement mean? (\textquotedblleft{}The dismissal, the layoff\textquotedblright{}.) Say it once more.
\end{tcolorbox}
\section[la démission]{la démission — the resignation}

\label{lesson:FR-C355-demission}



\begin{quote}
Before the new one: say the French for to dismiss, then the French for the dismissal.
\end{quote}

\subsection*{You'll want to know: la démission}
\textbf{la démission} — \textquotedblleft{}the resignation\textquotedblright{}.

\begin{quote}
Il a donné sa démission lundi.
\end{quote}

\begin{quote}
\emph{He handed in his resignation on Monday.}
\end{quote}

\emph{Donner sa démission} — to resign.

\begin{cousinweb}
\emph{démission} is a learned borrowing from Latin \emph{dēmissiō} \textquotedblleft{}a letting down, a lowering\textquotedblright{} (EtymDB; Lewis \& Short), from \emph{dēmittere} \textquotedblleft{}to send down\textquotedblright{}, built on \emph{mittere} \textquotedblleft{}to send\textquotedblright{}.
\end{cousinweb}

\begin{grammarlens}[title={What you've built}]
One more word for telling, explaining and describing at B1.
\end{grammarlens}

\subsection*{Your turn}
\begin{itemize}
\raggedright
  \item \emph{Say it:} \emph{la démission}
  \item \emph{Say it:} \emph{la démission}, once more
  \item \emph{Say it:} a sentence of your own with \emph{la démission}
  \item \emph{Recall:} say the French for to dismiss, then the French for the dismissal, then say \emph{la démission} again
\end{itemize}

\begin{tcolorbox}[breakable,colback=teal!4,colframe=teal!35!black,title={Before you move on}]
What does la démission mean? (\textquotedblleft{}The resignation\textquotedblright{}.) Say it once more.
\end{tcolorbox}
\section[démissionner]{démissionner — to resign, to quit (a job)}

\label{lesson:FR-C355-demissionner}



\begin{quote}
Before the new one: say the French for the dismissal, then the French for the resignation.
\end{quote}

\subsection*{You'll want to know: démissionner}
\textbf{démissionner} — \textquotedblleft{}to resign, to quit (a job)\textquotedblright{}.

\begin{quote}
Elle a démissionné parce qu'elle n'était pas heureuse.
\end{quote}

\begin{quote}
\emph{She resigned because she was not happy.}
\end{quote}

\emph{Démissionner de son poste}.

\begin{cousinweb}
\emph{démissionner} is built on \emph{démission}, from Latin \emph{dēmissiō} \textquotedblleft{}a letting down\textquotedblright{}, from \emph{dēmittere} \textquotedblleft{}to send down\textquotedblright{}, built on \emph{mittere} \textquotedblleft{}to send\textquotedblright{} (Lewis \& Short).
\end{cousinweb}

\begin{grammarlens}[title={What you've built}]
One more word for telling, explaining and describing at B1.
\end{grammarlens}

\subsection*{Your turn}
\begin{itemize}
\raggedright
  \item \emph{Say it:} \emph{démissionner}
  \item \emph{Say it:} \emph{démissionner}, once more
  \item \emph{Say it:} a sentence of your own with \emph{démissionner}
  \item \emph{Recall:} say the French for the dismissal, then the French for the resignation, then say \emph{démissionner} again
\end{itemize}

\begin{tcolorbox}[breakable,colback=teal!4,colframe=teal!35!black,title={Before you move on}]
What does démissionner mean? (\textquotedblleft{}To resign, to quit (a job)\textquotedblright{}.) Say it once more.
\end{tcolorbox}
\section[le préavis]{le préavis — the notice (period)}

\label{lesson:FR-C355-preavis}



\begin{quote}
Before the new one: say the French for the resignation, then the French for to resign.
\end{quote}

\subsection*{You'll want to know: le préavis}
\textbf{le préavis} — \textquotedblleft{}the notice (period)\textquotedblright{}.

\begin{quote}
Je dois faire un mois de préavis avant de partir.
\end{quote}

\begin{quote}
\emph{I have to work a month of notice before leaving.}
\end{quote}

\emph{Sans préavis} — without notice.

\begin{cousinweb}
\emph{préavis} is built from \emph{pré-} \textquotedblleft{}before\textquotedblright{} and \emph{avis} \textquotedblleft{}opinion, notice\textquotedblright{}. \emph{avis} comes from Old French \emph{avis}, built on \emph{vis} \textquotedblleft{}sight, view\textquotedblright{} (EtymDB), from Latin \emph{vīsus} \textquotedblleft{}a look, a sight\textquotedblright{}, from \emph{vidēre} \textquotedblleft{}to see\textquotedblright{} (Lewis \& Short).
\end{cousinweb}

\begin{grammarlens}[title={What you've built}]
One more word for telling, explaining and describing at B1.
\end{grammarlens}

\subsection*{Your turn}
\begin{itemize}
\raggedright
  \item \emph{Say it:} \emph{le préavis}
  \item \emph{Say it:} \emph{le préavis}, once more
  \item \emph{Say it:} a sentence of your own with \emph{le préavis}
  \item \emph{Recall:} say the French for the resignation, then the French for to resign, then say \emph{le préavis} again
\end{itemize}

\begin{tcolorbox}[breakable,colback=teal!4,colframe=teal!35!black,title={Before you move on}]
What does le préavis mean? (\textquotedblleft{}The notice (period)\textquotedblright{}.) Say it once more.
\end{tcolorbox}
